

\iffalse

\subsubsection{Application: Machine Unlearning for Trojan Model Cleanse}




%\vspace{2mm}
\noindent
\textbf{Setup.}
We demonstrate how machine unlearning can be applied to eliminate the effects of backdoor or poisoned data in a learned model. In this experiment, we implement the backdoor attack following the typical backdoor framework \citep{zeng2023narcissus,liu2022backdoor}. Specifically, an adversary modifies a small subset of training instances by embedding a backdoor trigger (e.g., a small image patch) and altering their labels toward a targeted incorrect class. The resulting Trojan model misclassifies any test input containing the trigger but behaves normally on data without the trigger.


We adopt machine unlearning as a defensive measure to remove the adverse influence of poisoned training examples on model predictions, following the rationale in \citep{liu2022backdoor}. We evaluate the performance of the unlearned model from two perspectives, backdoor attack success rate ({BSR}) and the TA on testing data. \Cref{application_of_backdoor_remove} shows BSR and TA of the Trojan model (with 100 poisoned samples) on CIFAR10.



\begin{figure}[t]
	\centering
	\vspace{-4mm}
	\subfloat{	\label{fig:1}  
		\includegraphics[scale=0.245]{../../../../PycharmProjects/Manifold_unl/Exp_results/On_CIFAR10/Backoor_clearance/CIFAR10_model_bac__mbu_r}
	}			
	\subfloat{ 		\label{fig:2}
		\includegraphics[scale=0.245]{../../../../PycharmProjects/Manifold_unl/Exp_results/On_CIFAR10/Backoor_clearance/CIFAR10_model_bac_TA__mbu_r}
	}
	\vspace{-2mm}
	\caption{Performance of Trojan model cleanse via proposed TCU on CIFAR10, where ``Origin'' refers to the original model with Trojan. \vspace{-0mm} } %\vspace{-2mm}
	\label{application_of_backdoor_remove} 
	\vspace{-4mm}
\end{figure}




\noindent
\textbf{Results.} As we can see in \Cref{application_of_backdoor_remove}, the original Trojan model has $100\%$ BSR, which means the model is fully backdoored. TCU can effectively reduce the BSR without too much accuracy degradation when the margin value $\alpha$ is less than 0.16. When $\alpha$ is larger than 0.16, although the backdoor cleanse effectiveness is much better, the model utility degradation will also be dramatic at the same time. In general, our proposed unlearning shows promise in the application of backdoor attack defense. Moreover, we should also notice that TCU achieves the backdoor cleanse without using the backdoor label information, which provides a wider application scenario such as clean-label backdooring removing \citep{zeng2023narcissus}.

\fi